% ===================== INTRODUCTION GÉNÉRALE =====================

Le marché immobilier tunisien reste largement structuré autour d'annonces
diffusées sur des sites généralistes, où la localisation d'un bien se résume le
plus souvent à un nom de ville ou de quartier saisi librement. Cette imprécision
rend la recherche fastidieuse pour l'acquéreur ou le locataire, qui ne peut ni
visualiser la répartition géographique de l'offre, ni filtrer autour d'un point
d'intérêt (lieu de travail, école, transport).

Le présent projet, réalisé au sein de \rapportEntreprise, vise à concevoir une
plateforme web dédiée à la \textbf{géolocalisation des biens immobiliers en
Tunisie}. Chaque bien y est positionné précisément sur une carte interactive,
rattaché au découpage administratif national (gouvernorat, délégation, localité)
et enrichi de coordonnées géographiques (latitude, longitude).

\paragraph{Problématique.}
\emph{Comment concevoir et modéliser une plateforme web permettant de publier,
géolocaliser et rechercher des biens immobiliers en Tunisie, de manière fiable et
ergonomique, en tenant compte des spécificités du découpage territorial
tunisien~?}

\paragraph{Objectif.}
Ce travail couvre exclusivement la \textbf{phase de conception}~: analyse du
contexte et des besoins, cahier des charges, modélisation UML (cas d'utilisation,
classes, séquence), conception de la base de données (MCD, MLD, dictionnaire de
données) et conception des interfaces sous Figma (arborescence, wireframes,
charte graphique, maquettes, prototype). La phase de développement n'entre pas
dans le périmètre de ce rapport.

\paragraph{Organisation du rapport.}
Le \textbf{chapitre~\ref{chap:contexte}} présente l'organisme d'accueil, l'étude
de l'existant, la problématique et le cahier des charges. Le
\textbf{chapitre~\ref{chap:solution}} introduit les concepts clés et la solution
proposée. Le \textbf{chapitre~\ref{chap:conception}} détaille la modélisation du
système et de la base de données. Le \textbf{chapitre~\ref{chap:uxui}} présente
la conception des interfaces utilisateur. Une conclusion dresse le bilan et
ouvre des perspectives.
